\documentclass[11pt,letterpaper]{article}
% Reusable style. Build each .tex from its own directory using scripts/build.sh.
\usepackage[T1]{fontenc}
\usepackage[utf8]{inputenc}
\usepackage{lmodern}
\usepackage[margin=0.86in,headheight=15pt]{geometry}
\usepackage{amsmath,amssymb,amsthm,mathtools}
\usepackage{booktabs,longtable,tabularx,array,multicol}
\usepackage[table]{xcolor}
\usepackage{microtype}
\usepackage{enumitem}
\usepackage{fancyhdr}
\usepackage{titlesec}
\usepackage{hyperref}
\usepackage{listings}
\usepackage[most]{tcolorbox}
\usepackage{needspace}
\definecolor{aeNavy}{HTML}{172D43}
\definecolor{aeBlue}{HTML}{176B90}
\definecolor{aeTeal}{HTML}{148477}
\definecolor{aePale}{HTML}{F2F7F9}
\definecolor{aeGray}{HTML}{5C6771}
\hypersetup{colorlinks=true,linkcolor=aeBlue,urlcolor=aeBlue,pdfauthor={Algorithm Engineering Documentation Suite},pdfsubject={Independent educational materials based on the supplied TAOCP volume topics}}
\titleformat{\section}{\Large\bfseries\color{aeNavy}}{\thesection}{0.6em}{}
\titleformat{\subsection}{\large\bfseries\color{aeBlue}}{\thesubsection}{0.6em}{}
\titleformat{\subsubsection}{\normalsize\bfseries\color{aeNavy}}{\thesubsubsection}{0.6em}{}
\setlist[itemize]{leftmargin=*,topsep=3pt,itemsep=1pt}
\setlist[enumerate]{leftmargin=*,topsep=3pt,itemsep=3pt}
\setlength{\parindent}{0pt}\setlength{\parskip}{5pt}
\pagestyle{fancy}\fancyhf{}\fancyhead[L]{\footnotesize\color{aeGray}ALGORITHM ENGINEERING}\fancyhead[R]{\footnotesize\color{aeGray}DOCUMENTATION SUITE}\fancyfoot[L]{\footnotesize\color{aeGray}Independent study and engineering practice}\fancyfoot[R]{\footnotesize\thepage}
\renewcommand{\headrulewidth}{0.35pt}
\newtcolorbox{engineeringnote}[1][]{breakable,colback=aePale,colframe=aeBlue,boxrule=0.5pt,arc=1mm,fonttitle=\bfseries,title={Engineering note},#1}
\newtcolorbox{evidencebox}[1][]{breakable,colback=white,colframe=aeTeal,boxrule=0.6pt,arc=1mm,fonttitle=\bfseries,title={Evidence and deliverables},#1}
\lstdefinestyle{algo}{basicstyle=\ttfamily\footnotesize,breaklines=true,frame=single,rulecolor=\color{aeGray!40},backgroundcolor=\color{aePale},showstringspaces=false,tabsize=2,columns=fullflexible,keepspaces=true,keywordstyle=\color{aeBlue}\bfseries,commentstyle=\color{aeTeal}}
\lstset{style=algo}
\newcommand{\DocTitle}[3]{%
\hypersetup{pageanchor=false}\begin{titlepage}\vspace*{1.2cm}
{\color{aeTeal}\rule{\linewidth}{2.3pt}}\par\vspace{1.15cm}
{\fontsize{24}{30}\selectfont\bfseries\color{aeNavy}\raggedright #1\par}\vspace{0.45cm}
{\Large\color{aeBlue}#2\par}\vspace{1.3cm}
{\large\bfseries #3\par}\vspace{0.35cm}
{\color{aeGray}Edition 1.0 \quad | \quad 29 October 2026\par}\vfill
{\color{aeGray}\small Original educational material aligned to the chapter and section topics in the five supplied volumes of Donald E. Knuth's \emph{The Art of Computer Programming}. This document is not an authorized reproduction or substitute for the original books.}\par\vspace{0.4cm}
{\color{aeTeal}\rule{\linewidth}{0.75pt}}
\end{titlepage}\hypersetup{pageanchor=true}\pagenumbering{roman}}
\newcommand{\SourceNote}[1]{\par\smallskip{\footnotesize\color{aeGray}\textbf{Source locator:} #1}\par\smallskip}
\newcommand{\Outcome}[1]{\par\textbf{Competency target.} #1\par}
\newcommand{\Checkpoint}[1]{\par\textbf{Checkpoint.} #1\par}
% Intentionally simple printable writing lines: call \AnswerLines{number}.
\newcount\alcount
\newcommand{\AnswerLines}[1]{\par\alcount=0\loop\ifnum\alcount<#1\par\noindent\rule{\linewidth}{0.3pt}\vspace{0.33cm}\advance\alcount by1\repeat\par}
\newcommand{\FormField}[2][0.62\linewidth]{\noindent\textbf{#2}\hspace{0.65em}\rule{#1}{0.4pt}\par\vspace{0.18cm}}
\newcommand{\ReviewTable}{\begin{tabularx}{\linewidth}{@{}p{0.23\linewidth}X p{0.2\linewidth}@{}}\toprule Criterion & Recorded evidence & Rating (0--4)\\\midrule Correctness & & \\[1em] Proof quality & & \\[1em] Complexity & & \\[1em] Testing & & \\[1em] Reproducibility & & \\[1em] Communication & & \\[1em]\bottomrule\end{tabularx}}

\begin{document}
\DocTitle{Five-Volume Source Crosswalk}{Primary sources, topic boundaries and engineering artifacts}{Engineering documentation}
\tableofcontents
\clearpage
\pagenumbering{arabic}
\section{Interpretation}
This index follows the supplied editions' tables of contents. Knuth's section numbers do not reset at physical volume boundaries. Volume 4B resumes Section 7.2.2 and is not Chapter 8. Asterisks on some source contents indicate specialized sections; this crosswalk does not assign a new mandatory sequence to those sections.
\small
\begin{longtable}{@{}p{0.14\linewidth}p{0.22\linewidth}p{0.31\linewidth}p{0.23\linewidth}@{}}
\toprule Source & Chapter/section & Primary focus & Engineering artifact\\\midrule\endfirsthead
\toprule Source & Chapter/section & Primary focus & Engineering artifact\\\midrule\endhead
Vol. 1 & 1.1--1.2.11 & Algorithms, induction, sums, asymptotics & Proof/complexity workbook\\
Vol. 1 & 1.3--1.4 & MIX, subroutines and programming techniques & Machine-model reflection\\
Vol. 1 & 2.1--2.2 & Information structures and linear lists & Queue/container lab\\
Vol. 1 & 2.3--2.6 & Trees, storage and allocation & Data structure review\\
Vol. 2 & 3.1--3.3 & Generators and statistical tests & PRNG experiment\\
Vol. 2 & 3.4--3.6 & Sampling, distributions, random sequences & Shuffle validation\\
Vol. 2 & 4.1--4.4 & Representation, floating point and precision & Arithmetic audit\\
Vol. 2 & 4.5--4.7 & GCD, polynomials and power series & Number theory lab\\
Vol. 3 & 5.1--5.3 & Permutations, sorting families, optimum sorting & Sorting suite\\
Vol. 3 & 5.4--5.5 & External sorting and trade-offs & I/O plan\\
Vol. 3 & 6.1--6.2 & Sequential and comparison searching & Search correctness proof\\
Vol. 3 & 6.3--6.5 & Digital searching, hashing and secondary keys & Index design assessment\\
Vol. 4A & 7.1.1--7.1.4 & Boolean basics, evaluation, bit tricks, BDDs & Bit-manipulation lab\\
Vol. 4A & 7.2.1.1--7.2.1.3 & Tuples, permutations, combinations & Iterator lab\\
Vol. 4A & 7.2.1.4--7.2.1.6 & Partitions and trees & Generation/counting workbook\\
Vol. 4B & 7.2.2 & Backtrack programming & Search-tree instrumentation\\
Vol. 4B & 7.2.2.1 & Dancing links, exact cover, colored covering & Exact-cover solver lab\\
Vol. 4B & 7.2.2.2 & SAT algorithms, clause learning and encoding & DPLL validation lab\\
\bottomrule
\end{longtable}
\normalsize
\section{Where to begin by outcome}
\textbf{Implement and optimize data structures:} Vol. 1 Chapter 2, then Vol. 3 Chapter 6, and only then specialized structures in Vol. 4B. \textbf{Numerical reliability:} Vol. 1 Section 1.2, then Vol. 2 Chapter 4. \textbf{Fast searching and indexing:} Vol. 1 Section 2.3, then Vol. 3 Sections 6.2--6.5. \textbf{Constraint solving:} Vol. 1 Section 1.2, Vol. 4A Sections 7.1 and 7.2.1, then Vol. 4B Sections 7.2.2.1--7.2.2.2.
\section{Source control and bibliographic provenance}
The source files are external references, not bundled into this deliverable. Keep the volumes in an authorized location; do not check copyright-protected copies into a public documentation repository. The accompanying README contains publication details; references in this suite identify volume and section rather than relying on fragile PDF page offsets.
\end{document}
